\Article{240}
\textsc{Theorem.} \textit{Suppose that from the forms $f,f'$ the form $F$ is composed; from $F$ and $f''$ the form $\mathfrak{F}$; from $f,f''$ the form $F'$; and from $F'$ and $f'$ the form $\mathfrak{F}'$. Then the forms $\mathfrak{F},\mathfrak{F}'$ are properly equivalent.}

\textit{Proof.} I. Let
\[
\begin{aligned}
f&=axx+2bxy+cyy\\
f'&=a'x'x'+2b'x'y'+c'y'y'\\
f''&=a''x''x''+2b''x''y''+c''y''y''\\
F&=AXX+2BXY+CYY\\
F'&=A'X'X'+2B'X'Y'+C'Y'Y'\\
\mathfrak{F}&=\mathfrak{A}\mathfrak{X}\mathfrak{X}+2\mathfrak{B}\mathfrak{X}\mathfrak{Y}+\mathfrak{C}\mathfrak{Y}\mathfrak{Y}\\
\mathfrak{F}'&=\mathfrak{A}'\mathfrak{X}'\mathfrak{X}'+2\mathfrak{B}'\mathfrak{X}'\mathfrak{Y}'+\mathfrak{C}'\mathfrak{Y}'\mathfrak{Y}'
\end{aligned}
\]
and let the determinants of these seven forms be $d,d',d'',D,D',\mathfrak{D},\mathfrak{D}'$, respectively. All have the same sign and the ratios of squares to one another. Further, let $m$ be the greatest common divisor of the numbers $a,2b,c$, and let $m',m'',M$ have similar meanings for the forms $f',f'';F$. Then, by the fourth conclusion of Article 235, $D$ is the greatest common divisor of the numbers $dm'm',d'mm$, and therefore $Dm''m''$ of the numbers $dm'm'm''m'',d'mmm''m''$; $M=mm'$; $\mathfrak{D}$ is the greatest common divisor of the numbers $Dm''m''$, $d''MM$, or of the numbers $Dm''m'',d''mmm'm'$. It follows that $\mathfrak{D}$ is the greatest common divisor of the three numbers $dm'm'm''m'',d'mmm''m'',d''mmm'm'$; but by a similar argument $\mathfrak{D}'$ is the greatest common divisor of the same three numbers. So, since $\mathfrak{D},\mathfrak{D}'$ have the same sign, we have $\mathfrak{D}=\mathfrak{D}'$, or the forms $\mathfrak{F},\mathfrak{F}'$ have the same determinant.

\SourceLeaf{261}
II. Now let $F$ pass into $ff'$ by the substitution
\[
\begin{aligned}
X&=pxx'+p'xy'+p''yx'+p'''yy'\\
Y&=qxx'+q'xy'+q''yx'+q'''yy'
\end{aligned}
\]
and $\mathfrak{F}$ into $Ff''$ by the substitution
\[
\begin{aligned}
\mathfrak{X}&=\mathfrak{p}Xx''+\mathfrak{p}'Xy''+\mathfrak{p}''Yx''+\mathfrak{p}'''Yy''\\
\mathfrak{Y}&=\mathfrak{q}Xx''+\mathfrak{q}'Xy''+\mathfrak{q}''Yx''+\mathfrak{q}'''Yy''
\end{aligned}
\]
and denote the positive square roots of $\frac{d}{D},\frac{d'}{D},\frac{D}{\mathfrak{D}}.\ \frac{d''}{\mathfrak{D}}$ by $n,n',\mathfrak{N},\mathfrak{n}''$. Then, by Article 235, we have eighteen equations, half belonging to the transformation of the form $F$ into $ff'$, the other half to that of the form $\mathfrak{F}$ into $Ff''$. The first is $pq'-qp'=an'$; the others can easily be formed on its model, and we omit them here for brevity. The quantities $n,n',\mathfrak{N},\mathfrak{n}''$, however, are rational, but not necessarily integers.

III. Substituting the values of $X,Y$ into those of $\mathfrak{X},\mathfrak{Y}$ gives a substitution of the following kind:
\[
\begin{aligned}
\mathfrak{X}={}&(1)xx'x''+(2)xx'y''+(3)xy'x''+(4)xy'y''\\
&+(5)yx'x''+(6)yx'y''+(7)yy'x''+(8)yy'y''\\
\mathfrak{Y}={}&(9)xx'x''+(10)xx'y''+(11)xy'x''+(12)xy'y''\\
&+(13)yx'x''+(14)yx'y''+(15)yy'x''+(16)yy'y''
\end{aligned}
\]
by which $\mathfrak{F}$ clearly passes into the product $ff'f''$. The coefficient $(1)$ is $=p\mathfrak{p}+q\mathfrak{p}''$; we do not give the values of the other fifteen, since anyone can work them out without difficulty. Denote the number $(1)(10)-(2)(9)$ by $(1,2)$, the number $(1)(11)-(3)(9)$ by $(1,3)$, and, in general, $(g)(8+h)-(h)(8+g)$ by $(g,h)$, supposing $g,h$ to be distinct integers between 1 and 16, the larger being $h$\HistoricalNote{*)}{The present meaning of these symbols must not be confused with the meaning they had in Article 234; for the numbers expressed by them \textit{here} correspond exactly to those which, in Article 234, were multiplied by the numbers denoted \textit{there} by similar symbols.}; in this way we have twenty-eight symbols in all. Now, denoting the positive square roots of $\frac{d}{\mathfrak{D}},\frac{d'}{\mathfrak{D}}$ by $\mathfrak{n},\mathfrak{n}'$ (which are $=n\mathfrak{N},n'\mathfrak{N}$), we derive the following 28 equations:

\SourceLeaf{262}
\[
\begin{aligned}
(1,2)&=aa'\mathfrak{n}''\\
(1,3)&=aa''\mathfrak{n}'\\
(1,4)&=ab'\mathfrak{n}''+ab''\mathfrak{n}'\\
(1,5)&=a'a''\mathfrak{n}\\
(1,6)&=a'b\mathfrak{n}''+a'b''\mathfrak{n}\\
(1,7)&=a''b\mathfrak{n}'+a''b'\mathfrak{n}\\
(1,8)&=bb'\mathfrak{n}''+bb''\mathfrak{n}'+b'b''\mathfrak{n}+\mathfrak{D}\mathfrak{n}\mathfrak{n}'\mathfrak{n}'',\\
(2,3)&=ab''\mathfrak{n}'-ab'\mathfrak{n}''\\
(2,4)&=ac''\mathfrak{n}'\\
(2,5)&=a'b''\mathfrak{n}-a'b\mathfrak{n}''\\
(2,6)&=a'c''\mathfrak{n}\\
(2,7)&=bb''\mathfrak{n}'+b'b''\mathfrak{n}-bb'\mathfrak{n}''-\mathfrak{D}\mathfrak{n}\mathfrak{n}'\mathfrak{n}'',\\
(2,8)&=bc''\mathfrak{n}'+b'c''\mathfrak{n}\\
(3,4)&=ac'\mathfrak{n}''
\end{aligned}
\]
\[
\begin{aligned}
(3,5)&=a''b'\mathfrak{n}-a''b\mathfrak{n}'\\
(3,6)&=bb'\mathfrak{n}''+b'b''\mathfrak{n}-bb''\mathfrak{n}'-\mathfrak{D}\mathfrak{n}\mathfrak{n}'\mathfrak{n}''\\
(3,7)&=a''c'\mathfrak{n}\\
(3,8)&=bc'\mathfrak{n}''+b''c\mathfrak{n}\\
(4,5)&=b'b''\mathfrak{n}-bb'\mathfrak{n}''-bb''\mathfrak{n}'+\mathfrak{D}\mathfrak{n}\mathfrak{n}'\mathfrak{n}''\\
(4,6)&=b'c''\mathfrak{n}-bc''\mathfrak{n}'\\
(4,7)&=b''c'\mathfrak{n}-bc'\mathfrak{n}''\\
(4,8)&=c'c''\mathfrak{n}\\
(5,6)&=ca'\mathfrak{n}''\\
(5,7)&=ca''\mathfrak{n}'\\
(5,8)&=b'c\mathfrak{n}''+b''c\mathfrak{n}'\\
(6,7)&=b''c\mathfrak{n}'-b'c\mathfrak{n}''\\
(6,8)&=cc''\mathfrak{n}'\\
(7,8)&=cc'\mathfrak{n}''
\end{aligned}
\]
which we will denote by $\Phi$, and nine others:
\[
\begin{aligned}
(10)(11)-(9)(12)&=a\mathfrak{n}'\mathfrak{n}''\mathfrak{A}\\
(1)(12)-(2)(11)-(3)(10)+(4)(9)&=2a\mathfrak{n}'\mathfrak{n}''\mathfrak{B}\\
(2)(3)-(1)(4)&=a\mathfrak{n}'\mathfrak{n}''\mathfrak{C}\\
-(9)(16)+(10)(15)+(11)(14)-(12)(13)&=2b\mathfrak{n}'\mathfrak{n}''\mathfrak{A}\\
\begin{gathered}
(1)(16)-(2)(15)-(3)(14)+(4)(13)\{\\
+(5)(12)-(6)(11)-(7)(10)+(8)(9)\}
\end{gathered}&=4b\mathfrak{n}'\mathfrak{n}''\mathfrak{B}\\
-(1)(8)+(2)(7)+(3)(6)-(4)(5)&=2b\mathfrak{n}'\mathfrak{n}''\mathfrak{C}\\
(14)(15)-(13)(16)&=c\mathfrak{n}'\mathfrak{n}''\mathfrak{A}\\
(5)(16)-(6)(15)-(7)(14)+(8)(13)&=2c\mathfrak{n}'\mathfrak{n}''\mathfrak{B}\\
(6)(7)-(5)(8)&=c\mathfrak{n}'\mathfrak{n}''\mathfrak{C}
\end{aligned}
\]
which we will denote by $\Psi$\HistoricalNote{*)}{We should note that 18 other equations like these $\Psi$ can be derived, with the factors $a$, $2b$, $c$ on the right replaced by $a'$, $2b'$, $c'$; $a''$, $2b''$, $c''$. But we omit them, since they are not needed for our purpose.}.

IV. Deriving the origin of all these 37 equations would take too long. It will be enough to verify a few, on whose model the others can be proved without much difficulty.

\leavevmode\SourceLeaf{263}
1) We have
\[
\begin{aligned}
(1,2)&=(1)(10)-(2)(9)\\
&=(\mathfrak{p}\mathfrak{q}'-\mathfrak{q}\mathfrak{p}')pp\\&\quad+(\mathfrak{p}\mathfrak{q}'''-\mathfrak{q}\mathfrak{p}'''-\mathfrak{p}'\mathfrak{q}''+\mathfrak{q}'\mathfrak{p}'')pq\\&\quad+(\mathfrak{p}''\mathfrak{q}'''-\mathfrak{q}''\mathfrak{p}''')qq\\
&=\mathfrak{n}''(App+2Bpq+Cqq)=\mathfrak{n}''aa'
\end{aligned}
\]
which is the first equation.

2) We obtain
\[
\begin{aligned}(1,3)&=(1)(11)-(3)(9)\\&=(\mathfrak{p}\mathfrak{q}''-\mathfrak{q}\mathfrak{p}'')(pq'-qp')\\&=a''\mathfrak{N}an'=aa''\mathfrak{n}'\end{aligned}
\]
the second equation.

3) We have
\[
\begin{aligned}
(1,8)&=(1)(16)-(8)(9)\\
&=(\mathfrak{p}\mathfrak{q}'-\mathfrak{q}\mathfrak{p}')pp'''\\&\quad+(\mathfrak{p}\mathfrak{q}'''-\mathfrak{q}\mathfrak{p}''')pq'''\\&\quad-(\mathfrak{p}'\mathfrak{q}''-\mathfrak{q}'\mathfrak{p}'')qp'''\\&\quad+(\mathfrak{p}''\mathfrak{q}'''-\mathfrak{q}''\mathfrak{p}''')qq'''\\
&=\mathfrak{n}''(App'''+B(pq'''+qp''')+Cqq''')\\&\quad+b''\mathfrak{N}(pq'''-qp''')\\
&=\mathfrak{n}''(bb'+\surd dd')+b''\mathfrak{N}(b'n+bn')\text{\textsuperscript{*)}}\\
&=\mathfrak{n}''bb'+\mathfrak{n}'bb''+\mathfrak{n}b'b''+\mathfrak{D}\mathfrak{n}\mathfrak{n}'\mathfrak{n}''
\end{aligned}
\]
\HistoricalDisplayNote{*)}{This follows from equation 10 of Article 235 and what follows it. The radical quantity $\surd dd'$ becomes $=Dnn'=\mathfrak{D}nn'\mathfrak{N}\mathfrak{N}=\mathfrak{D}\mathfrak{n}\mathfrak{n}'$.}
the eighth equation in $\Phi$. We leave the remaining equations for the readers to verify.

V. It follows from equations $\Phi$ that the twenty-eight numbers $(1,2)$, $(1,3)$, etc. have no common divisor, as follows. First we observe that the twenty-seven products of three factors, whose first factor is either $\mathfrak{n}$, second one of the numbers $a'$, $2b'$, $c'$, and third one of the numbers $a''$, $2b''$, $c''$; or whose first is $\mathfrak{n}'$, second one of the numbers $a$, $2b$, $c$, and third one of the numbers $a''$, $2b''$, $c''$; or, finally, whose first is $\mathfrak{n}''$, second one of the numbers $a$, $2b$, $c$, and third one of the numbers $a'$, $2b'$, $c'$---each of these twenty-seven products, by equations $\Phi$, equals either one of the twenty-eight numbers $(1,2)$, $(1,3)$, etc., or a sum or difference of several of them (for \textit{example}, $\mathfrak{n}a'a''=(1,5)$, $2\mathfrak{n}a'b''=(1,6)+(2,5)$, $4\mathfrak{n}b'b''=(1,8)+(2,7)+(3,6)+(4,5)$, and so on). So if these numbers had a common divisor, it would necessarily divide all those products as well. From this, with the help of Article 40 and the method often used above, we readily deduce that the same divisor must also divide the numbers $\mathfrak{n}m'm''$, $\mathfrak{n}'mm''$, $\mathfrak{n}''mm'$, and therefore their
\SourceLeaf{264}
squares, which are $\frac{dm'm'm''m''}{\mathfrak{D}}$, $\frac{d'mmm''m''}{\mathfrak{D}}$, $\frac{d''mmm'm'}{\mathfrak{D}}$, would be divisible by its square, \textit{which is absurd}, since by I the greatest common divisor of the three numerators is $\mathfrak{D}$, and therefore the squares themselves can have no common divisor.

VI. All this concerns the transformation of the form $\mathfrak{F}$ into $ff'f''$, and was derived from the transformations of the form $F$ into $ff'$ and of the form $\mathfrak{F}$ into $Ff''$. But in exactly the same way, from the transformations of the form $F'$ into $ff''$ and of the form $\mathfrak{F}'$ into $F'f'$, we derive a transformation of the form $\mathfrak{F}'$ into $ff'f''$ of the following kind:
\[
\begin{aligned}
\mathfrak{X}'&=(1)'xx'x''+(2)'xx'y''+(3)'xy'x''+\text{etc.}\\
\mathfrak{Y}'&=(9)'xx'x''+(10)'xx'y''+(11)'xy'x''+\text{etc.}
\end{aligned}
\]
(denoting all coefficients as in the transformation of the form $\mathfrak{F}$ into $ff'f''$, and adding a prime to each to distinguish them). From this, as before, we derive 28 equations analogous to $\Phi$, which we will denote by $\Phi'$, and nine others analogous to $\Psi$, which we will denote by $\Psi'$. Namely, denoting
\[
(1)'(10)'-(2)'(9)'\text{ by }(1,2)',\quad
(1)'(11)'-(3)'(9)'\text{ by }(1,3)'
\]
the equations $\Phi'$ are
\[
(1,2)'=aa'\mathfrak{n}'',\quad (1,3)'=aa''\mathfrak{n}',\quad\text{etc.}
\]
and the equations $\Psi'$ are
\[
(10)'(11)'-(9)'(12)'=a\mathfrak{n}'\mathfrak{n}''\mathfrak{A}',\quad\text{etc.}
\]
(For brevity we leave the fuller working to the readers; those familiar with the subject will find that a new calculation is not even necessary, and that the first analysis is easily carried over by analogy.) With this done, from $\Phi$ and $\Phi'$ we immediately have
\[
(1,2)=(1,2)',\quad (1,3)=(1,3)',\quad (1,4)=(1,4)',\quad (2,3)=(2,3)',\quad\text{etc.}
\]
From this, and because all of $(1,2)$, $(1,3)$, $(2,3)$, etc. have no common divisor (by V), the lemma of Article 234 shows that four integers $\alpha$, $\beta$, $\gamma$, $\delta$ can be determined such that
\[
\begin{aligned}
\alpha(1)'+\beta(9)'&=(1),\\\alpha(2)'+\beta(10)'&=(2),\\\alpha(3)'+\beta(11)'&=(3),\quad\text{etc.}\\
\gamma(1)'+\delta(9)'&=(9),\\\gamma(2)'+\delta(10)'&=(10),\\\gamma(3)'+\delta(11)'&=(11),\quad\text{etc.}
\end{aligned}
\]
and $\alpha\delta-\beta\gamma=1$.

\SourceLeaf{265}
VII. From this, substituting from the first three equations $\Psi$ the values of $a\mathfrak{A}$, $a\mathfrak{B}$, $a\mathfrak{C}$, and from the first three equations $\Psi'$ the values of $a\mathfrak{A}'$, $a\mathfrak{B}'$, $a\mathfrak{C}'$, we readily verify that
\[
\begin{aligned}
a(\mathfrak{A}\alpha\alpha+2\mathfrak{B}\alpha\gamma+\mathfrak{C}\gamma\gamma)&=a\mathfrak{A}'\\
a(\mathfrak{A}\alpha\beta+\mathfrak{B}(\alpha\delta+\beta\gamma)+\mathfrak{C}\gamma\delta)&=a\mathfrak{B}'\\
a(\mathfrak{A}\beta\beta+2\mathfrak{B}\beta\delta+\mathfrak{C}\delta\delta)&=a\mathfrak{C}'
\end{aligned}
\]
Unless $a=0$, it clearly follows that the form $\mathfrak{F}$ passes into $\mathfrak{F}'$ by the proper substitution $\alpha$, $\beta$, $\gamma$, $\delta$. --- But using the next three equations in place of the first three in $\Psi$ and $\Psi'$, we readily verify three equations entirely analogous to those just given, with the factor $a$ everywhere replaced by $b$. So the same conclusion clearly still holds, provided we do not have $b=0$. Finally, using the last three equations $\Psi$, $\Psi'$, we find in the same way that the conclusion is true unless $c=0$. Therefore, since certainly not all of $a$, $b$, $c$ can be $=0$ at once, the form $\mathfrak{F}$ necessarily passes by the substitution $\alpha$, $\beta$, $\gamma$, $\delta$ into $\mathfrak{F}'$, and so is properly equivalent to this form. \textit{Q. E. D.}

\Needspace{8\baselineskip}
\Article{241}
We will call a form such as $\mathfrak{F}$ or $\mathfrak{F}'$, obtained by composing one of three given forms with the form composed of the other two, \textit{composed of these three forms}. The preceding article shows that the order in which the three forms are composed makes no difference. Similarly, given any number of forms $f$, $f'$, $f''$, $f'''$, etc. (whose determinants must have the ratios of squares to one another), if we compose the form $f$ with $f'$, the resulting form with $f''$, the next result with $f'''$, etc., the form obtained at the end of this operation is said to be composed \textit{of all the forms} $f$, $f'$, $f''$, $f'''$, etc. It is easy to prove that here too the order of composition is arbitrary; \textit{i.e.}, whatever the order in which these forms are composed, the forms resulting from the composition are always properly equivalent. --- Further, if the forms $f$, $f'$, $f''$, etc. are properly equivalent to the forms $g$, $g'$, $g''$, etc., respectively, then the form composed of the latter is clearly properly equivalent to that composed of the former.
